% GENERATED FILE. Edit canonical lessons, then run npm run generate:books.
% canonical-source-hash: fnv1a64:28cdaed4ab3c52a4
% canonical-lessons: MW-C233-palan, MW-C233-nakel, MW-C233-saj, MW-C233-kasno, MW-C233-ghughro

\chapter{Camel Gear}
\label{ch:mw-camel-gear}
\hlchaptermodality{233}

\begin{chapteropening}
\textbf{By the end of this chapter:} \emph{I can name a camel's saddle, rein, trappings and bells.}

Complete the last lesson of chapter 233: I can name a camel's saddle, rein, trappings and bells.
\end{chapteropening}

\section[\texorpdfstring{\mw{पलाण} (palāṇ)}{पलाण (palāṇ)}]{\mw{पलाण} (palāṇ) — a camel saddle; a pack-saddle}

\label{lesson:MW-C233-palan}



\begin{quote}
Before the new one: say the Marwadi for a pebble, then the Marwadi for a jujube bush.
\end{quote}

\subsection*{You'll want to know: \mw{पलाण}}
\textbf{\mw{पलाण}} — \emph{palāṇ} — \textquotedblleft{}a camel saddle; a pack-saddle\textquotedblright{}.

\begin{quote}
\mw{ऊंट} \mw{माथै} \mw{नुवो} \mw{पलाण} \mw{है।}
\end{quote}

\begin{quote}
\emph{ū̃ṭ māthai nuvo palāṇ hai.}
\end{quote}

\begin{quote}
The camel has a new saddle.
\end{quote}

A masculine noun.

\begin{cousinweb}
\textbf{\mw{पलाण}} (\emph{palāṇ}) is inherited from Sanskrit \emph{palyāṇa}, \textquotedblleft{}a saddle\textquotedblright{}; CDIAL lists Hindi \emph{palān} and Gujarati \emph{palāṇ}, \textquotedblleft{}a pack-saddle\textquotedblright{}. CDIAL gives no Marwari form in this entry (CDIAL 7966).
\end{cousinweb}

\begin{grammarlens}[title={What you've built}]
One more word for this chapter.
\end{grammarlens}

\subsection*{Your turn}
\begin{itemize}
\raggedright
  \item \emph{Say it:} \emph{palāṇ}
  \item \emph{Say it:} \emph{palāṇ}, once more
  \item \emph{Say it:} \emph{ū̃ṭ māthai nuvo palāṇ hai.} — The camel has a new saddle.
  \item \emph{Recall:} say the Marwadi for a pebble, then the Marwadi for a jujube bush, then say \emph{palāṇ} again
\end{itemize}

\begin{tcolorbox}[breakable,colback=teal!4,colframe=teal!35!black,title={Before you move on}]
What does \mw{पलाण} mean? (\textquotedblleft{}A camel saddle; a pack-saddle\textquotedblright{}.) Say it once more.
\end{tcolorbox}
\section[\texorpdfstring{\mw{नकेल} (nakel)}{नकेल (nakel)}]{\mw{नकेल} (nakel) — a camel's nose-peg and rein}

\label{lesson:MW-C233-nakel}



\begin{quote}
Before the new one: say the Marwadi for a jujube bush, then the Marwadi for a camel saddle.
\end{quote}

\subsection*{You'll want to know: \mw{नकेल}}
\textbf{\mw{नकेल}} — \emph{nakel} — \textquotedblleft{}a camel's nose-peg and rein\textquotedblright{}.

\begin{quote}
\mw{नकेल} \mw{ढीली} \mw{है।}
\end{quote}

\begin{quote}
\emph{nakel ḍhīlī hai.}
\end{quote}

\begin{quote}
The nose-rein is loose.
\end{quote}

A feminine noun.

\begin{cousinweb}
\textbf{\mw{नकेल}} (\emph{nakel}) is built from a reconstructed \emph{*nakka}, \textquotedblleft{}nose\textquotedblright{}, and Sanskrit \emph{kīla}, \textquotedblleft{}a peg\textquotedblright{}; CDIAL lists Hindi \emph{nakel}, the pin fixed in a camel's nose. CDIAL gives no Marwari form in this entry (CDIAL 6910).
\end{cousinweb}

\begin{grammarlens}[title={What you've built}]
One more word for this chapter.
\end{grammarlens}

\subsection*{Your turn}
\begin{itemize}
\raggedright
  \item \emph{Say it:} \emph{nakel}
  \item \emph{Say it:} \emph{nakel}, once more
  \item \emph{Say it:} \emph{nakel ḍhīlī hai.} — The nose-rein is loose.
  \item \emph{Recall:} say the Marwadi for a jujube bush, then the Marwadi for a camel saddle, then say \emph{nakel} again
\end{itemize}

\begin{tcolorbox}[breakable,colback=teal!4,colframe=teal!35!black,title={Before you move on}]
What does \mw{नकेल} mean? (\textquotedblleft{}A camel's nose-peg and rein\textquotedblright{}.) Say it once more.
\end{tcolorbox}
\section[\texorpdfstring{\mw{साज} (sāj)}{साज (sāj)}]{\mw{साज} (sāj) — trappings (of a camel or a horse); harness}

\label{lesson:MW-C233-saj}



\begin{quote}
Before the new one: say the Marwadi for a camel saddle, then the Marwadi for a camel's nose-peg.
\end{quote}

\subsection*{You'll want to know: \mw{साज}}
\textbf{\mw{साज}} — \emph{sāj} — \textquotedblleft{}trappings (of a camel or a horse); harness\textquotedblright{}.

\begin{quote}
\mw{ऊंट} \mw{माथै} \mw{लाल} \mw{साज} \mw{है।}
\end{quote}

\begin{quote}
\emph{ū̃ṭ māthai lāl sāj hai.}
\end{quote}

\begin{quote}
The camel wears red trappings.
\end{quote}

A masculine noun.

\begin{cousinweb}
\textbf{\mw{साज}} (\emph{sāj}) is inherited from Sanskrit \emph{sajja}, \textquotedblleft{}equipped, put in place\textquotedblright{}; CDIAL lists Old Marwari \emph{sāja}, \textquotedblleft{}accoutrements\textquotedblright{} (CDIAL 13095).
\end{cousinweb}

\begin{grammarlens}[title={What you've built}]
One more word for this chapter.
\end{grammarlens}

\subsection*{Your turn}
\begin{itemize}
\raggedright
  \item \emph{Say it:} \emph{sāj}
  \item \emph{Say it:} \emph{sāj}, once more
  \item \emph{Say it:} \emph{ū̃ṭ māthai lāl sāj hai.} — The camel wears red trappings.
  \item \emph{Recall:} say the Marwadi for a camel saddle, then the Marwadi for a camel's nose-peg, then say \emph{sāj} again
\end{itemize}

\begin{tcolorbox}[breakable,colback=teal!4,colframe=teal!35!black,title={Before you move on}]
What does \mw{साज} mean? (\textquotedblleft{}Trappings (of a camel or a horse); harness\textquotedblright{}.) Say it once more.
\end{tcolorbox}
\section[\texorpdfstring{\mw{कसणो} (kasṇo)}{कसणो (kasṇo)}]{\mw{कसणो} (kasṇo) — to tighten; to harness}

\label{lesson:MW-C233-kasno}



\begin{quote}
Before the new one: say the Marwadi for a camel's nose-peg, then the Marwadi for trappings (of a camel or a horse).
\end{quote}

\subsection*{You'll want to know: \mw{कसणो}}
\textbf{\mw{कसणो}} — \emph{kasṇo} — \textquotedblleft{}to tighten; to harness\textquotedblright{}.

\begin{quote}
\mw{हूं} \mw{ऊंट} \mw{माथै} \mw{पलाण} \mw{कसूं} \mw{हूं।}
\end{quote}

\begin{quote}
\emph{hū̃ ū̃ṭ māthai palāṇ kasū̃ hū̃.}
\end{quote}

\begin{quote}
I tighten the saddle on the camel.
\end{quote}

Past (I …): \mw{म्हैं} \mw{पलाण} \mw{कस्यो}, \emph{mhaĩ palāṇ kasyo}.

\begin{cousinweb}
\textbf{\mw{कसणो}} (\emph{kasṇo}) is inherited from Sanskrit \emph{karṣati}, \textquotedblleft{}draws, pulls\textquotedblright{}; CDIAL lists the Old Marwari verb \emph{kasaï}, \textquotedblleft{}harnesses, binds\textquotedblright{}, and Hindi \emph{kasnā}, \textquotedblleft{}to tighten\textquotedblright{} (CDIAL 2908).
\end{cousinweb}

\begin{grammarlens}[title={What you've built}]
One more word for this chapter.
\end{grammarlens}

\subsection*{Your turn}
\begin{itemize}
\raggedright
  \item \emph{Say it:} \emph{kasṇo}
  \item \emph{Say it:} \emph{kasṇo}, once more
  \item \emph{Say it:} \emph{hū̃ ū̃ṭ māthai palāṇ kasū̃ hū̃.} — I tighten the saddle on the camel.
  \item \emph{Recall:} say the Marwadi for a camel's nose-peg, then the Marwadi for trappings (of a camel or a horse), then say \emph{kasṇo} again
\end{itemize}

\begin{tcolorbox}[breakable,colback=teal!4,colframe=teal!35!black,title={Before you move on}]
What does \mw{कसणो} mean? (\textquotedblleft{}To tighten; to harness\textquotedblright{}.) Say it once more.
\end{tcolorbox}
\section[\texorpdfstring{\mw{घूघरो} (ghūghro)}{घूघरो (ghūghro)}]{\mw{घूघरो} (ghūghro) — a small bell (on a harness or an anklet)}

\label{lesson:MW-C233-ghughro}



\begin{quote}
Before the new one: say the Marwadi for trappings (of a camel or a horse), then the Marwadi for to tighten.
\end{quote}

\subsection*{You'll want to know: \mw{घूघरो}}
\textbf{\mw{घूघरो}} — \emph{ghūghro} — \textquotedblleft{}a small bell (on a harness or an anklet)\textquotedblright{}.

\begin{quote}
\mw{हूं} \mw{घूघरो} \mw{सुणूं} \mw{हूं।}
\end{quote}

\begin{quote}
\emph{hū̃ ghūghro suṇū̃ hū̃.}
\end{quote}

\begin{quote}
I hear the little bell.
\end{quote}

A masculine noun.

\begin{cousinweb}
\textbf{\mw{घूघरो}} (\emph{ghūghro}) is inherited from a reconstructed \emph{*ghuṅghura}, \textquotedblleft{}a bell\textquotedblright{}, which CDIAL takes to imitate the sound; CDIAL lists Old Marwari \emph{ghūghara}, \textquotedblleft{}a string of bells\textquotedblright{} (CDIAL 4477).
\end{cousinweb}

\begin{grammarlens}[title={What you've built}]
One more word for this chapter.
\end{grammarlens}

\subsection*{Your turn}
\begin{itemize}
\raggedright
  \item \emph{Say it:} \emph{ghūghro}
  \item \emph{Say it:} \emph{ghūghro}, once more
  \item \emph{Say it:} \emph{hū̃ ghūghro suṇū̃ hū̃.} — I hear the little bell.
  \item \emph{Recall:} say the Marwadi for trappings (of a camel or a horse), then the Marwadi for to tighten, then say \emph{ghūghro} again
\end{itemize}

\begin{tcolorbox}[breakable,colback=teal!4,colframe=teal!35!black,title={Before you move on}]
What does \mw{घूघरो} mean? (\textquotedblleft{}A small bell (on a harness or an anklet)\textquotedblright{}.) Say it once more.
\end{tcolorbox}
